\section{One iteration of $K$-means}
\label{sec:ex1}
\begin{problem}
One-dimensional observations are $X=\{1,3,9,12\}$, with $K=2$ and initial centroids $\mu_1^{(0)}=1$, $\mu_2^{(0)}=12$.
\begin{enumerate}[label=\textbf{(\alph*)}]
\item Assign each observation to its nearest centroid using Euclidean distance.
\item Compute the updated centroids $\mu_1^{(1)}$ and $\mu_2^{(1)}$.
\item Calculate inertia $\sum_i\norm{x_i-\mu_{c(i)}}_2^2$ using the \emph{updated} centroids.
\item If you repeat assignment with those updated centroids, does any cluster membership change? Explain.
\end{enumerate}
\end{problem}


\begin{concept}
$K$-means alternates two operations: (i) assign each point to its nearest centroid; (ii) replace each centroid with its assigned points' arithmetic mean. Its objective is the \emph{within-cluster sum of squares} $J=\sum_i\norm{x_i-\mu_{c(i)}}_2^2$.
\end{concept}

\subq{(a) Assign each point to its closest initial centroid}
In one dimension, Euclidean distance is absolute difference: $d(x,\mu)=\abs{x-\mu}$. Compute both distances before choosing the smaller value:
\begin{center}
\begin{tabular}{@{}rrrrl@{}}\toprule
$x_i$ & $\abs{x_i-1}$ & $\abs{x_i-12}$ & Nearest & Assignment\\\midrule
1&0&11&1&$C_1$\\
3&2&9&1&$C_1$\\
9&8&3&2&$C_2$\\
12&11&0&2&$C_2$\\\bottomrule
\end{tabular}
\end{center}
Thus $C_1=\{1,3\}$ and $C_2=\{9,12\}$.

\subq{(b) Update the centroids}
Take the mean of each assigned cluster:
\[
\mu_1^{(1)}=\frac{1+3}{2}=\boxed{2},\qquad
\mu_2^{(1)}=\frac{9+12}{2}=\boxed{10.5}.
\]
The updated centroids minimize the within-cluster squared-error objective \emph{for the current assignments}.

\subq{(c) Calculate the inertia using the \emph{updated} centroids}
Substitute each point's assigned, updated centroid into the specified objective:
\begin{align*}
J^{(1)}&=(1-2)^2+(3-2)^2+(9-10.5)^2+(12-10.5)^2\\
&=1+1+2.25+2.25=\boxed{6.5}.
\end{align*}
Note that this is a \emph{sum}, not a mean; do not divide by four.

\subq{(d) Reassign using the updated centroids}
The nearest-centroid decision boundary in one dimension is their midpoint:
\[
b=\frac{2+10.5}{2}=6.25.
\]
Values $1$ and $3$ are less than $6.25$; $9$ and $12$ exceed it. Consequently the assignments remain $\{1,3\}$ and $\{9,12\}$. A subsequent centroid update also gives $2$ and $10.5$: this solution has reached a fixed point of the alternating algorithm.
\begin{answer}
\textbf{Result:} $C_1=\{1,3\}$, $C_2=\{9,12\}$; centroids $(2,10.5)$; inertia $6.5$; no further assignment changes.
\end{answer}
